\chapter{Extraits de code commentés}
\label{anx:code}

\section{Implémentation de RICE}
\begin{lstlisting}[caption={Scoring RICE normalisé sur [0;100].}]
def score_rice(f: Feature) -> float:
    """RICE = (Reach * Impact * Confidence) / Effort."""
    raw = (f.reach * f.impact * f.confidence) / max(f.effort, 0.1)
    return min(100.0, raw * 1.5)  # normalisation empirique
\end{lstlisting}

\section{Implémentation de WSJF}
\begin{lstlisting}[caption={Scoring WSJF (SAFe).}]
def score_wsjf(f: Feature) -> float:
    """WSJF = Cost of Delay / Job Size."""
    cod = f.business_value + f.time_criticality + f.risk_reduction
    raw = cod / max(f.job_size, 0.1)
    return min(100.0, raw * 4.0)
\end{lstlisting}

\section{Implémentation d'ICE}
\begin{lstlisting}[caption={Scoring ICE.}]
def score_ice(f: Feature) -> float:
    """ICE = Impact * Confidence * Ease, normalise sur [0;100]."""
    raw = f.impact * f.confidence * f.ease
    return min(100.0, raw * 2.5)
\end{lstlisting}

\section{Sprint Planner glouton}
\begin{lstlisting}[caption={Algorithme glouton du Sprint Planner.}]
def plan_sprint(features, velocity, algorithm):
    scored = prioritize(features, algorithm)
    selected, deferred = [], []
    cumulative = 0.0
    for rank, item in enumerate(scored, start=1):
        if cumulative + item.effort <= velocity:
            selected.append({"rank": rank, **item.dict()})
            cumulative += item.effort
        else:
            deferred.append({"rank": rank, **item.dict()})
    utilization = round(100 * cumulative / velocity, 1)
    return {"selected": selected, "deferred": deferred,
            "utilization": utilization, "velocity": velocity}
\end{lstlisting}

\section{Construction de la chaîne RAG (LangChain)}
\begin{lstlisting}[caption={Pipeline RAG avec mémoire conversationnelle.}]
def build_rag_chain(model="gpt-4o", session_id="default"):
    llm         = get_llm(model)
    vectorstore = get_vectorstore()
    retriever   = vectorstore.as_retriever(
        search_type="mmr",
        search_kwargs={"k": 6, "fetch_k": 20},
    )
    memory = ConversationBufferWindowMemory(
        k=10, memory_key="chat_history",
        return_messages=True, output_key="answer",
    )
    return ConversationalRetrievalChain.from_llm(
        llm=llm, retriever=retriever, memory=memory,
        combine_docs_chain_kwargs={"prompt": PO_PROMPT},
        return_source_documents=True,
    )
\end{lstlisting}
